% Part: computability
% Chapter: computability-theory
% Section: non-comp-set

\documentclass[../../../include/open-logic-section]{subfiles}

\begin{document}

\olfileid{cmp}{thy}{ncp}
\olsection{અસંગણનીય ગણો અસ્તિત્વમાં છે}


ઉપર જોયું કે દરેક સંગણનીય ગણ સંગણકીય રીતે પરિગણનીય છે. શું
વિરુદ્ધ દાવો પણ સાચો છે? આગળનું પ્રમેય બતાવે છે કે સામાન્ય રીતે
તે સાચો નથી.

\begin{thm}\ollabel{thm:K-0}
$K_0$ એ ગણ $\Setabs{\tuple{e, x}}{\cfind{e}(x) \fdefined}$ હોય.
તો $K_0$ સંગણકીય રીતે પરિગણનીય છે, પરંતુ સંગણનીય નથી.
\end{thm}

\begin{proof}
$K_0$ સંગણકીય રીતે પરિગણનીય છે તે જોવા, નોંધો કે તે નીચેના
વિધેય~$f$નો પ્રદેશ છે:
\[
f(z) = \umin{y}{(\len{z} = 2 \land T((z)_0, (z)_1, y))}.
\]
કારણ કે $\cfind{e}(x)$ વ્યાખ્યાયિત હોય તો
$f(\tuple{e, x})$ થંભતી સંગણનાનો અનુક્રમ શોધે છે; જો
$\cfind{e}(x)$ અનિર્ધારિત હોય, તો $f(\tuple{e, x})$ પણ
અનિર્ધારિત છે; અને $z$ જોડીનો કોડ પણ ન હોય, તો $f(z)$ પણ
અનિર્ધારિત છે.

$K_0$ સંગણનીય નથી એ હકીકત અટકવાની સમસ્યાની અનિર્ણેયતા,
\olref[hlt]{thm:halting-problem}, જ છે.
\end{proof}

ગણ $K_0$ એ એવી જોડીઓ $\tuple{e,x}$નો ગણ છે કે
$\cfind{e}(x) \fdefined$; એટલે $\tuple{e,x} \in K_0$ ત્યારે
અને માત્ર ત્યારે જ્યારે $\cfind{e}$ ઇનપુટ~$x$ પર વ્યાખ્યાયિત
હોય, અર્થાત્ અટકે. તેથી તેને ``અટકવાનો ગણ'' પણ કહે છે.
ગણ $K = \Setabs{e}{\cfind{e}(e) \fdefined}$ એ
``સ્વ-અટકવાનો ગણ'' છે. તેને ઘણી વાર પ્રમાણભૂત અનિર્ણેય ગણ
તરીકે વાપરાય છે.

\begin{thm}\ollabel{thm:K}
સ્વ-અટકવાનો ગણ $K = \Setabs{e}{\cfind{e}(e) \fdefined}$
!!{c.e.} છે, પરંતુ નિર્ણેય નથી.
\end{thm}

\begin{proof}
  માનો કે $K$ નિર્ણેય છે, એટલે તેનું લક્ષણ વિધેય
  $\Char{K}$ સંગણનીય છે. માનો કે
  \[d(e) = \begin{cases}
  1 & \text{જો\/ $\Char{K}(e) = 0$}\\
  \fundefined & \text{અન્યથા.}
  \end{cases}
  \]
  $k$, $d$નો સૂચક હોય, એટલે $d \simeq \cfind{k}$.
  તો $d(k) \simeq \cfind{k}(k)$. પરંતુ $d$ની વ્યાખ્યા મુજબ
  $d(k) \fdefined$ ત્યારે અને માત્ર ત્યારે જ્યારે
  $\cfind{k}(k) \fundefined$; આ વિરોધાભાસ છે.

  $K$, $f(x) = \umin{y}{T(x,x,y)}$નો પ્રદેશ છે અને તેથી
  !!{c.e.} છે.
\end{proof}

\end{document}
